% Source: https://artofproblemsolving.com/community/c3h2555488_triangle_problem
% Stable identity: aops-topic-2555488-q01
We have $QA=QP$ and $QR=QB$, and $\angle AQP=\angle RQB=60^\circ$.

\begin{asy}
usepackage("amsmath");
size(6cm);
defaultpen(fontsize(10pt));
pair Q=(0,0),P=(-3,-1.2),R=(2.6,-1.6);
pair outer(pair X,pair Y,pair O){pair c1=rotate(60,X)*Y,c2=rotate(-60,X)*Y;
 return (cross(Y-X,c1-X)*cross(Y-X,O-X)<0)?c1:c2;}
pair A=outer(Q,P,R),B=outer(Q,R,P);
draw(P--Q--R--cycle,linewidth(0.8));
draw(A--P--Q--cycle);draw(B--Q--R--cycle);
draw(A--R,dashed+blue);draw(P--B,dashed+red);
label("$P$",P,SW);label("$Q$",Q,S);label("$R$",R,SE);label("$A$",A,NW);label("$B$",B,NE);
\end{asy}

Consider the rotation about $Q$ through $60^\circ$ that takes $A$ to $P$ (possible because $QA=QP$ and $\angle AQP=60^\circ$). Since both equilateral triangles are drawn outside $PQR$, the turn from ray $QA$ to ray $QP$ and the turn from ray $QR$ to ray $QB$ are in the same direction: going around $Q$ we meet the rays in the order $QA$, $QP$, $QR$, $QB$, with $60^\circ$ from $QA$ to $QP$ and $60^\circ$ from $QR$ to $QB$. So the same rotation takes $R$ to $B$ (as $QR=QB$).

A rotation preserves distances, and it takes the segment $AR$ to the segment $PB$. Hence $AR=PB$. $\blacksquare$

(Equivalently, triangles $AQR$ and $PQB$ are congruent by SAS: $QA=QP$, $QR=QB$, and the angles $AQR$ and $PQB$ are both obtained by adding $60^\circ$ to $\angle PQR$, or both equal $300^\circ-\angle PQR$ if that sum exceeds $180^\circ$.) The hypothesis that the triangle is obtuse with longest side $PR$ is not needed.
